\begin{theorem}[余项 \texorpdfstring{$B$}{B} 比特的二阶重整化恒等式与 InfoNCE 缩放律]\label{thm:pom-infonce-remainder-bits-second-order-renormalization}
沿用推论 \ref{cor:pom-infonce-fold-ledger-remainder-bits} 的 \(g_B(\omega)=(X,R_B)\) 与纤维占用数 \(d_{m,B}(x,r)\)。
定义余项分区的二阶碰撞和与其归一化版本
\[
S_{2,B}(m):=\sum_{x\in X_m}\sum_{r\in\{0,1\}^B} d_{m,B}(x,r)^2,
\qquad
M_{2,B}(m):=\frac{S_{2,B}(m)}{4^m}=\sum_{x,r}\Bigl(\frac{d_{m,B}(x,r)}{2^m}\Bigr)^2.
\]
对每个 \(x\in X_m\) 作 Euclid 分解
\[
d_m(x)=2^B q_x+s_x,\qquad q_x=\Big\lfloor\frac{d_m(x)}{2^B}\Big\rfloor,\quad s_x=d_m(x)\bmod 2^B\in\{0,\dots,2^B-1\}.
\]
则有精确恒等式
\[
\boxed{\;
S_{2,B}(m)=2^{-B}S_2(m)-2^{-B}\sum_{x\in X_m}s_x^2+\sum_{x\in X_m}s_x\;}
\]
并因此得到二阶碰撞的重整化渐近
\[
M_{2,B}(m)=2^{-B}M_2(m)+O\!\left(\frac{2^B|X_m|}{4^m}\right),
\qquad
\frac{M_{2,B}(m)}{M_2(m)}\longrightarrow 2^{-B},
\]
其中 \(M_2(m)=S_2(m)/4^m\)。
\par
进一步，对固定整数 \(K\ge 2\) 定义余项分区的 Bayes 最优 InfoNCE 值
\(
L_{K,B}^\star(m):=\inf_s L_{K,B}(s)
\)
（见推论 \ref{cor:pom-infonce-fold-ledger-remainder-bits}）。
则当 \(m\to\infty\) 时有
\[
L_{K,B}^\star(m)
=2^{-B}(K-1)\log 2\cdot M_2(m)\bigl(1+o(1)\bigr),
\qquad
\frac{L_{K,B}^\star(m)}{L_{K}^\star(m)}\longrightarrow 2^{-B},
\]
其中 \(L_K^\star(m)\) 为定理 \ref{thm:pom-bayes-infonce-second-collision-dominance} 的 Bayes 最优值。
\end{theorem}
\begin{proof}
由推论 \ref{cor:pom-infonce-fold-ledger-remainder-bits}，对每个固定 \(x\)，模 \(2^B\) 的占用数满足：恰有 \(s_x\) 个余类取值为 \(q_x+1\)，其余 \(2^B-s_x\) 个余类取值为 \(q_x\)。
因此
\[
\sum_{r}d_{m,B}(x,r)^2
=s_x(q_x+1)^2+(2^B-s_x)q_x^2
=2^Bq_x^2+2s_xq_x+s_x.
\]
另一方面，
\(
d_m(x)^2=(2^Bq_x+s_x)^2=2^{2B}q_x^2+2^{B+1}q_xs_x+s_x^2
\)，
从而逐点恒等式
\[
2^{-B}d_m(x)^2-2^{-B}s_x^2+s_x
=2^Bq_x^2+2s_xq_x+s_x
=\sum_{r}d_{m,B}(x,r)^2
\]
成立。对 \(x\) 求和得到所示的 \(S_{2,B}\) 重整化恒等式。
\par
记修正项
\(
\Delta_B(m):=\sum_x\bigl(s_x-2^{-B}s_x^2\bigr)
\)。
由于 \(0\le s_x<2^B\)，并且函数 \(s\mapsto s-2^{-B}s^2\) 在区间 \([0,2^B]\) 上有界（例如绝对值 \(O(2^B)\)），故
\(
|\Delta_B(m)|=O(2^B|X_m|)
\)。
两边除以 \(4^m\) 即得
\(
M_{2,B}(m)=2^{-B}M_2(m)+O(2^B|X_m|/4^m)
\)。
又因 \(|X_m|=F_{m+2}=\Theta(\varphi^m)\) 且 \(M_2(m)=\Theta((r_2/4)^m)\)（其中 \(r_2>\varphi\)，见例如推论 \ref{cor:pom-s2-asymptotic} 的数值与谱证书），故误差项相对 \(M_2(m)\) 衰减至零，从而 \(M_{2,B}(m)/M_2(m)\to 2^{-B}\)。
\par
最后，对分区 \(g_B\) 重复定理 \ref{thm:pom-bayes-infonce-second-collision-dominance} 的二项近似论证，可得
\(
L_{K,B}^\star(m)=(K-1)\log 2\cdot M_{2,B}(m)\bigl(1+o(1)\bigr)
\)；
与 \(M_{2,B}(m)\sim 2^{-B}M_2(m)\) 合并即得所示缩放律与比值收敛。证毕。
\end{proof}

\begin{corollary}[InfoNCE 的指数不变量与 \texorpdfstring{$r_2$}{r2} 的一致可识别性]\label{cor:pom-infonce-exponential-invariant-r2-identifiability}
在定理 \ref{thm:pom-infonce-remainder-bits-second-order-renormalization} 的设定下，对任意固定整数 \(K\ge 2\)、\(B\ge 0\)，极限存在且
\[
\lim_{m\to\infty}\bigl(L_{K,B}^\star(m)\bigr)^{1/m}=\frac{r_2}{4}.
\]
并且 \(r_2\) 可由任意一组固定 \((K,B)\) 的 Bayes 最优 InfoNCE 序列一致反演：
\[
\boxed{\;
r_2
=4\lim_{m\to\infty}\left(\frac{2^B\,L_{K,B}^\star(m)}{(K-1)\log 2}\right)^{1/m}\; }.
\]
\end{corollary}
\begin{proof}
由定理 \ref{thm:pom-infonce-remainder-bits-second-order-renormalization}，
\(
L_{K,B}^\star(m)
=2^{-B}(K-1)\log 2\cdot M_2(m)\,(1+o(1))
\)。
对两边取 \(m\)-次根并用 \(M_2(m)=S_2(m)/4^m=\Theta((r_2/4)^m)\) 得到第一式；
第二式为同一渐近的等价重排。证毕。
\end{proof}

\paragraph{InfoNCE 的有限分辨率层析重建}
\begin{theorem}[由 Bayes 最优 InfoNCE 总损失重建纤维谱]\label{thm:pom-infonce-fiber-spectrum-tomography}
固定分辨率 \(m\)，令 \(N:=|X_m|\)，并以 \(w_m(x)=d_m(x)/2^m\) 记折叠推前分布。
对每个 \(K\ge 2\) 记
\(
L_K^\star(m):=\inf_s L_K(s)
\)
为定理 \ref{thm:pom-infonce-partition-optimal-closed-form} 在 \(g=\Fold_m\) 下的 Bayes 最优值。
则多重集 \(\{w_m(x)\}_{x\in X_m}\)（等价地，\(\{d_m(x)\}_{x\in X_m}\)）可由有限数据 \(\{L_K^\star(m)\}_{K=2}^{N}\) 唯一重建。
\par
更具体地，令
\[
\Psi_K(p):=\EE\bigl[\log(1+J)\bigr],\qquad J\sim\mathrm{Binomial}(K-1,p),
\]
并令 \(\Phi_K(p):=p\,\Psi_K(p)\)。
则 \(\Phi_K\) 是次数 \(K\) 的多项式且 \(\Phi_K(0)=\Phi_K'(0)=0\)，从而存在常数 \(\{\beta_{K,q}\}_{2\le q\le K}\) 使得
\[
\Phi_K(p)=\sum_{q=2}^{K}\beta_{K,q}\,p^q.
\]
定义幂和矩
\(
M_q(m):=\sum_{x\in X_m}w_m(x)^q
\)。
则
\[
L_K^\star(m)=\sum_{q=2}^{K}\beta_{K,q}\,M_q(m)\qquad(2\le K\le N),
\]
并且系数矩阵 \((\beta_{K,q})_{2\le K\le N,\,2\le q\le K}\) 为上三角且对角元
\[
\beta_{K,K}
=\sum_{j=0}^{K-1}(-1)^{K-1-j}\binom{K-1}{j}\log(1+j)
=\Delta^{K-1}\bigl(\log(1+\cdot)\bigr)(0)\ \neq\ 0,
\]
因此可由上三角回代从 \(\{L_K^\star(m)\}_{K=2}^{N}\) 解出 \(\{M_q(m)\}_{q=2}^{N}\)。
由 Newton 恒等式可进一步从幂和 \(\{M_q\}_{q=1}^{N}\)（其中 \(M_1=1\)）恢复多项式
\[
P_m(t):=\prod_{x\in X_m}\bigl(t-w_m(x)\bigr),
\]
从而以根的多重集形式重建 \(\{w_m(x)\}\)，并经 \(d_m(x)=2^m w_m(x)\) 重建纤维谱。
\end{theorem}
\begin{proof}
由定理 \ref{thm:pom-infonce-partition-optimal-closed-form} 在 \(g=\Fold_m\) 下，
\[
L_K^\star(m)=\sum_{x\in X_m}w_m(x)\,\Psi_K\!\bigl(w_m(x)\bigr)
=\sum_{x\in X_m}\Phi_K\!\bigl(w_m(x)\bigr).
\]
由于 \(\Psi_K\) 是二项分布在有限支持 \(\{0,\dots,K-1\}\) 上的期望，且其概率质量函数是 \(p\) 的多项式，故 \(\Psi_K\) 为次数 \(K-1\) 的多项式并满足 \(\Psi_K(0)=0\)；
于是 \(\Phi_K(p)=p\Psi_K(p)\) 为次数 \(K\) 的多项式且在 \(0\) 处至少二重零点，从而存在 \(\beta_{K,q}\) 使 \(\Phi_K(p)=\sum_{q=2}^{K}\beta_{K,q}p^q\)。
代回 \(p=w_m(x)\) 并对 \(x\) 求和得到
\(
L_K^\star(m)=\sum_{q=2}^{K}\beta_{K,q}M_q(m)
\)。
\par
为说明可逆性，注意 \(\beta_{K,K}\) 为 \(\Phi_K\) 的最高次系数，即 \(\Psi_K\) 中 \(p^{K-1}\) 的系数。
由
\[
\Psi_K(p)=\sum_{j=0}^{K-1}\log(1+j)\binom{K-1}{j}p^{j}(1-p)^{K-1-j}
\]
展开 \((1-p)^{K-1-j}\) 并收集 \(p^{K-1}\) 项，即得所示交错和表达式，它等于 \(\log(1+\cdot)\) 在 \(0\) 处的 \((K-1)\) 阶前向差分 \(\Delta^{K-1}(\log(1+\cdot))(0)\)。
由于 \(\log(1+\cdot)\) 非多项式，故其所有高阶差分均不恒为零，因而 \(\beta_{K,K}\neq 0\)。
于是 \(\{L_K^\star(m)\}_{K=2}^{N}\) 与 \(\{M_q(m)\}_{q=2}^{N}\) 之间给出上三角、对角非零的线性系统，可由回代唯一解出 \(\{M_q(m)\}_{q=2}^{N}\)。
\par
最后，已知幂和 \(M_1,\dots,M_N\) 可由 Newton 恒等式恢复 \(P_m\) 的初等对称系数，从而恢复根的多重集 \(\{w_m(x)\}\)；再由 \(d_m(x)=2^m w_m(x)\) 得到纤维谱。证毕。
\end{proof}

\endinput
